% Generated draft: requires mathematical and source review.
\documentclass[11pt]{article}
\usepackage[margin=25mm]{geometry}
\usepackage{fontspec}
\setmainfont{Latin Modern Roman}
\usepackage{amsmath,amssymb,mathtools,unicode-math}
\setmathfont{Latin Modern Math}
\usepackage{enumitem,needspace,xurl,hyperref}
\setlength{\parindent}{0pt}
\setlength{\parskip}{5pt}
\setlength{\emergencystretch}{3em}
\newcommand{\N}{\mathbb{N}}
\newcommand{\Z}{\mathbb{Z}}
\newcommand{\Q}{\mathbb{Q}}
\newcommand{\R}{\mathbb{R}}
\newcommand{\C}{\mathbb{C}}
\newcommand{\F}{\mathbb{F}}
\newcommand{\A}{\mathbb{A}}
\newcommand{\PP}{\mathbb{P}}
\newcommand{\E}{\mathbb{E}}
\newcommand{\eps}{\varepsilon}
\newcommand{\dd}{\,\mathrm{d}}
\begin{document}
\section*{UnsolvedMath ID: 30004200}
\textbf{Source problem code:} catalogue entry
\subsection{Definitions and mathematical statement}
The displayed source statement does not specify the class of groups \(G\), the meaning of \(\operatorname{Sub}(G)\), or whether \(0\in\mathbb N\). Its background discusses compact Lie groups. The following formulation uses that interpretation and makes the standard chromatic conventions explicit; these choices require confirmation against the cited report.

Fix a prime \(p\) and a compact Lie group \(G\). Let \(\operatorname{Sub}(G)\) denote the set of closed subgroups of \(G\), and write \(\mathbb Z_{\geq 0}=\{0,1,2,\ldots\}\). Work with genuine \(G\)-spectra indexed on a complete \(G\)-universe. A finite \(p\)-local \(G\)-spectrum means a retract, in the \(p\)-local equivariant stable homotopy category, of a finite \(G\)-cell spectrum. For a closed subgroup \(H\leq G\), let \(\Phi^H X\) denote the underlying nonequivariant geometric fixed-point spectrum of the restriction of \(X\) to \(H\).

For \(n\geq 1\), let \(K(n)\) denote Morava \(K\)-theory of height \(n\) at \(p\), and set \(K(0)=H\mathbb Q\). For a finite \(p\)-local nonequivariant spectrum \(Y\), define
\[
\operatorname{type}_p(Y)=
\begin{cases}
\min\{n\in\mathbb Z_{\geq 0}:K(n)_*(Y)\neq 0\},&\text{if this set is nonempty},\\
\infty,&\text{otherwise}.
\end{cases}
\]
The chromatic type function of \(X\) is
\[
 t_X\colon\operatorname{Sub}(G)\longrightarrow\mathbb Z_{\geq 0}\cup\{\infty\},
 \qquad t_X(H)=\operatorname{type}_p(\Phi^H X).
\]
This function is constant on conjugacy classes. If the source instead uses \(\operatorname{Sub}(G)\) for conjugacy classes of closed subgroups, the same definition descends to that quotient.

The realization problem is to characterize, for each such \(G\) and \(p\), exactly those functions
\[
 f\colon\operatorname{Sub}(G)\longrightarrow\mathbb Z_{\geq 0}\cup\{\infty\}
\]
for which
\[
 \exists\,X\text{ a finite genuine }p\text{-local }G\text{-spectrum}\quad
 \forall\,H\leq G\text{ closed},\qquad
 f(H)=\operatorname{type}_p(\Phi^H X).
\]
The supplied statement imposes no explicit continuity hypothesis on \(f\); any necessary regularity is part of the requested characterization. The compact Lie hypothesis, genuine equivariant interpretation, geometric fixed-point convention, and inclusion of height zero are explicit interpretive choices because the supplied statement does not define them.
\subsection{Short English statement}
For a fixed prime \(p\) and compact Lie group \(G\), determine which assignments of chromatic heights to closed subgroups arise from the geometric fixed points of one finite \(p\)-local genuine \(G\)-spectrum. The conventions concerning \(G\), subgroups, and height zero require source confirmation. Supplied dated curation labels the problem partially solved, reporting a classification for abelian compact Lie groups and an unresolved full nonabelian classification; this assessment has not been independently verified.
\subsection{Sources}
\begin{itemize}
\item Supplied local UnsolvedMath registry record, permanent ID \(30004200\), title \emph{Realizable Type Functions of Equivariant Spectra}, source problem code \texttt{catalogue entry}, in the Oberwolfach Reports Open Problems collection. No individual proposer is supplied; the proposed year is recorded as 2019.
\item Supplied report citation: \emph{Homotopy Theory (2020)}, Oberwolfach Reports, pp.\ 2183--2256, DOI \texttt{10.4171/owr/2019/36}. Supplied source URL: \url{https://doi.org/10.4171/owr/2019/36}. The supplied literature list also links the original report as \url{https://doi.org/10.4171/OWR/2019/36}.
\item Supplied literature source: \emph{Balmer Spectrum for Finite Abelian Groups}, \url{https://arxiv.org/abs/1709.04828}.
\item Supplied literature source: \emph{Balmer Spectrum for Compact Lie Groups}, \url{https://arxiv.org/abs/1810.04698}.
\item Literature assessment checked at \texttt{2026-08-22}: the supplied curation reports that the type-function realization problem is classified for abelian compact Lie groups, while the full nonabelian classification remains unresolved. This is supplied dated curation, not an independently verified assessment. The record expressly requests independent verification of time-sensitive status and cited claims.
\item The supplied statement status is \texttt{corrected\_verified}, with an accompanying claim that mathematical content was preserved while notation, typography, and sentence structure were repaired. This is registry metadata, not an independent verification of the present formulation.
\end{itemize}
\end{document}
